\documentclass[11pt]{article}
\usepackage{amsmath,amssymb,amsthm,mathtools,hyperref}

\begin{document}
\section*{Independent origins of multicellularity}
\label{prob:307}
\label{rank:199}
\par\smallskip

\noindent\textbf{Domain.} Multicellularity; repeated transitions. \quad \textbf{Retained status/source context.} Open explanation of lineage-dependent transition rates and differentiation. Experimental transitions are established, so the existence of evolvable multicellular organization is not a blanket open claim.

\paragraph{Definitions and assumptions.} On a dated species phylogeny $\mathcal T$, code independently defined states $0$ (unicellular), $1$ (cohesive groups) and $2$ (persistent differentiated multicellularity). State 2 requires declared reproductive integration and cell-type criteria; aggregation alone is insufficient. A continuous-time Markov model has off-diagonal rates $q_{ij}(x)=\exp(\beta_{ij}^{\top}x)$ per million years for prespecified allowed transitions $i\to j$, with $q_{ij}(x)=0$ for forbidden off-diagonal transitions and $q_{ii}(x)=-\sum_{j\ne i}q_{ij}(x)$. The ecological/genomic covariates $x$ are prespecified; their dimension, scaling and coefficient parameterization must make $\beta_{ij}^{\top}x$ dimensionless. Trait misclassification, uncertain ancestral branches, extinct lineages and nonrandom taxon sampling have a stated observation model. Origins are counted as qualifying state transitions, not merely changes in body size.

\paragraph{Problem statement.} Estimate which covariates predict transition to state2 and test transfer:
    \[\lambda_2(t,x)=\sum_{i<2}\mathbb{P}(Z_t=i\mid Z_t<2,x)q_{i2}(x).\]
    Determine identifiable contrasts in $\lambda_2(t,\cdot)$ between declared at-risk lineage classes, with confidence or compatible-model bounds, and predict independently held-out transitions. A fitted rate is not automatically a causal ecological effect.

\paragraph{Additional binding assumptions and resolution criterion.} Here $Z_t\in\{0,1,2\}$ denotes the coded state of a lineage at time or branch position $t$ under the instance's declared lineage-selection and evolutionary law. That law must specify the time coordinate, initial/root law, branching inheritance and dependence, and treatment of extinct and unsampled lineages; whether the tree is conditioned on or jointly modeled remains an instance declaration. A rate of entry to state 2 among at-risk lineages conditions on $Z_t<2$ and depends on the lineage-state distribution at a stated time or branch position. The displayed conditional weights are defined where $\mathbb{P}(Z_t<2\mid x)>0$; for continuously distributed $x$, they use a declared version of the conditional law on its stated domain. Population-wide transition incidence uses different unnormalized weights. Specify whether re-entry after loss counts as a new independent origin or a recurrence retaining ancestral machinery. The displayed $\lambda_2$ is an entry intensity and is not automatically a first-independent-origin hazard when previously differentiated lineages can return to the at-risk set; the origin-counting choice remains to be declared. Time-varying covariates and hidden trait states need a joint evolutionary law; a regression coefficient alone is not causal evidence.

\paragraph{Alternative formulations.} (i) A comparative reconstruction estimates origin counts conditional on coding and phylogeny. (ii) A replicated experimental lineage model measures transition probabilities per generation in a restricted setting. (iii) A cell-group mechanics model predicts size or cohesion, which is weaker than stable differentiation and reproductive integration. Connecting these scales requires an explicit mapping and timescale assumptions.

\paragraph{Obstructions and partial progress.} Long-term yeast experiments demonstrate macroscopic multicellular evolution. An October 2026 review synthesizes further transition mechanisms; neither establishes one rate law for all lineages. Sparse independent origins and correlated traits limit comparative causal inference. Lost transitions, reversible grouping and threshold-dependent classification change apparent rates. The residual target is lineage-specific, validated prediction of integrated complexity while separating existing experimental achievements from universal historical explanations.

\paragraph{Sources (supplied provenance; not all rechecked).} \begin{itemize}
\item West et al., Major evolutionary transitions in individuality (2015). \url{https://pmc.ncbi.nlm.nih.gov/articles/PMC4547252/} (Indexed review abstract checked via PubMed 25964342; direct PMC open encountered a browser challenge.)
\item Bozdag et al., De novo evolution of macroscopic multicellularity (2023). \url{https://pubmed.ncbi.nlm.nih.gov/37165189/} (Primary experiment checked; a restricted evolutionary achievement.)
\item Overcoming key challenges in the evolutionary transition to multicellularity (2026). \url{https://www.nature.com/articles/s41586-026-11038-w} (Indexed review checked, published 7 October 2026; direct full-page retrieval failed.)
\item Indexed abstract for the supplied West et al. review. \url{https://pubmed.ncbi.nlm.nih.gov/25964342/}
\end{itemize}


\section*{Binding catalogue conventions}
This edition revises both supplied catalogues and preserves original problem numbers 1--330. It adds one globally unique label \verb|\label{rank:x}| per problem, where smaller $x$ means greater provisional difficulty. A rank label is an anchor at that problem's section: \verb|\ref{rank:x}| returns the section number, not the literal difficulty number. The visible \texttt{rank:x} text and the index give the difficulty number itself.

The entries are \emph{parameterized mathematical research schemes}, including empirical prediction, causal identification, historical reconstruction, theory and engineering. They are not 330 certified unsolved mathematical conjectures. A rigorous instance fixes the objects and restrictions below; changing those objects changes the question. Established identities, finite calculations, non-identification constructions and solved special cases are not being posed anew as open. Some residual tasks concern informative bounds, biological realization or external validation rather than a new theorem.

Every main definition and problem statement received a mathematical editorial review. Each entry now has an additional binding-conditions paragraph; direct corrections are separately recorded in the accompanying audit. The alternative formulations inherit the applicable conventions, but are different targets unless equivalence is proved. Where an old broad phrase conflicts with an explicit restriction below it, the explicit restriction controls.

\textbf{Literature boundary.} Thirteen selective primary-literature checks cover eighteen entries in this revision; their access levels and scope are recorded explicitly. The supplied bibliography and its earlier source-access descriptions are retained as provenance, not represented as newly verified. Neither this review nor a source's description proves the absence of later solutions. The date of this edition is 13 October 2026; there is no claim of exhaustive current-status certification for all 330 entries.

\subsection*{Binding conventions for a concrete instance}
\paragraph{Objects, quantifiers and non-vacuity.}
An instance declares measurable state, input, output and parameter spaces; a fixed candidate family; admissible initial laws, interventions and environments; a finite horizon or specified asymptotic regime; observation and sampling laws; and performance metrics, positive normalization scales and positive tolerances. These are chosen before fitting or testing, not adjusted after seeing failures. Shared parameters must genuinely be shared across the declared validation domain. Finite dimension alone does not restrict arbitrary real-number encodings: specify a finite library, bounded description length, or regularity/complexity bounds. Exclude constant or no-signal solutions through a fixed informative-output or baseline-improvement requirement where appropriate. A universally quantified result cannot be established by choosing one trivial favorable instance.

\paragraph{Probability laws and data.}
Every expectation or probability refers to the specified stochastic law, including initial uncertainty, process noise, observation noise and sampling as appropriate. A deterministic equation is not a probability model without a declared random input or initial law. ``Observed'' targets in population-level formulas mean population observables $g_j^*$, not noisy estimates asserted to be exact. Let $\widetilde g_j$ be independent validation estimates and suppose
\[
\mathbb{P}\{\lvert\widetilde g_j-g_j^*\rvert\le r_j\text{ for every }j\}\ge1-\alpha.
\]
Then an uncertainty-aware sufficient certification rule for positive target tolerances $\varepsilon_j$ is
\[
\max_j\frac{|\widehat g_j-\widetilde g_j|+r_j}{\varepsilon_j}\le1.
\]
On that simultaneous confidence event it implies every population error is at most its tolerance, by the triangle inequality. Vector/norm versions use corresponding valid simultaneous bounds. Training uncertainty must also be included when the claim averages over repeated training data. Evaluation on finitely many conditions does not establish a continuous-domain supremum without covering and regularity bounds.

\paragraph{Dynamics and well-posedness.}
For ODEs and SDEs give coefficient domains, initial laws and existence/uniqueness conditions in the intended solution sense, or keep the entire admitted solution set. For stochastic jump models specify nonnegative off-diagonal rates, diagonal normalization, rate orientation and non-explosion. For PDEs specify spatial domains, boundary/initial conditions, conservation or source terms, and a weak or classical solution class. Moving-domain concentration equations include dilution/transport terms. A stationary law is used only for a time-homogeneous augmented model under existence assumptions; a stationary approximation to a driven system must be justified. Absorbing extinction states are retained. Survival-conditioned, quasi-stationary and unconditional laws are distinct and must not be substituted for one another.

\paragraph{Identification and causal meaning.}
Two parameter values are observationally equivalent when they induce the same full distribution under all allowed observation/intervention designs, not merely the same likelihood value for one sample. A target is identified only if constant on the equivalence class. Gauge symmetries, label switching and reference conventions are quotiented or fixed. Finite-data compatible sets use a declared confidence, likelihood or discrepancy rule; an empty set signals model rejection. Bayesian credible sets are not automatically confidence sets. A prior or regularizer does not create data identification. Bounds intended to resolve a question must be informative at a predeclared width or decision threshold.

A causal contrast uses well-defined intervention versions and a stated causal model. Random assignment, or explicit consistency, overlap and exchangeability assumptions, is required for causal interpretation from observations. Interference and time-varying confounding are represented when present. Fixing a mediator by intervention differs from conditioning on its observed value. Marginal potential-outcome laws do not determine individual cross-world effects or an arbitrary joint coupling. Ethically admissible designs, rather than unconstrained interventions, define the experimental menu.

\paragraph{Metrics, units, optimization and limits.}
Specify the sampling or reference measure for densities and integrals, the log base, and dimensionless arguments of logarithms. Total variation is $\sup_A|P(A)-Q(A)|$; Wasserstein distances require the declared ground metric and appropriate moment conditions; KL divergence uses its stated absolute-continuity convention. All normed readouts are scaled consistently, and error probabilities use $0<\alpha<1$. A ratio has a positive denominator or an explicit undefined/zero convention. Thermodynamic entropy production is nonnegative in expectation under the relevant hypotheses; individual stochastic realizations need not be. An infimum is not an attained minimum without compactness/coercivity and continuity or another existence argument. Fix the order of large-population, long-time and vanishing-mutation limits. Stability claims declare perturbation space and symmetries.

\subsection*{How the provisional difficulty order was assigned}
Four editorial components are scored from 1 (lower burden) to 5 (higher burden): (a) formal or inferential difficulty, (b) access to evidence and identification burden, (c) mechanistic/computational integration, and (d) independent validation and transfer. Their unweighted total orders the entries. For selected entries in the highest bands, an explicit editorial priority breaks equal totals; all other ties are broken by original problem number. Equal totals are \emph{not} evidence of reliably different true difficulties. Scores are judgments, not measurements, completion-time estimates or consensus rankings.

The object ranked is an informative, nondegenerate resolution of the \emph{scoped} scientific target, with the stated budget and validation domain. A sharp useful impossibility result can qualify; an already known identity, vacuous bound, arbitrary favorable model or lookup table cannot. Rank 330 is merely the least difficult relative assessment in this catalogue, not a claim of an easy biological solution. Broad historical titles do not automatically outrank precisely formulated mechanistic tasks.



\end{document}
